\chapter{Background and Related Work}
\label{ch:background}

\section{Static Analysis for Vulnerability Detection}

The oldest automated approach to vulnerability detection is static analysis: examining source
code or compiled artefacts without executing them, using rules derived from expert knowledge
of common vulnerability patterns.
Tools such as Flawfinder~\citep{wheeler2004flawfinder} and RATS scan C and C++ source code
for calls to known-dangerous functions --- \texttt{strcpy}, \texttt{gets}, \texttt{sprintf} ---
and flag their usage.
Commercial tools such as Coverity and Fortify extend this with dataflow analysis that tracks
whether values from untrusted sources can reach dangerous sinks, reducing some categories of
false positives.

Despite their maturity, rule-based static analysers have three well-documented limitations.
First, their rules must be written and maintained by security experts, making them expensive
to extend to new vulnerability classes or new languages.
Second, they cannot represent complex semantic conditions --- a \texttt{strcpy} call is flagged
regardless of whether the destination buffer is large enough, because tracking buffer sizes
through arbitrary control flow is undecidable in general.
Third, false-positive rates are high: industry surveys suggest that between 50\% and 90\%
of alerts from commercial tools are false positives, which leads to alert fatigue and
analyst disengagement~\citep{johnson2013don}.

Machine learning offers a complementary approach: rather than encoding rules, learn from
labelled examples, with the potential to generalise to patterns that rules do not capture.

\section{Early Machine Learning Approaches}

Early ML-based vulnerability detection treated source functions as sequences of tokens or
program slices and applied recurrent neural networks.
VulDeePecker~\citep{li2018vuldeepecker} extracted program slices centred on API calls
and classified them with a bidirectional LSTM, achieving strong results on two CWE classes
(buffer overflows and resource management errors) in C/C++.
SySeVR~\citep{li2022sysevr} extended this to a broader range of vulnerability types by
defining syntax-based, semantics-based, vector representations of program slices.

These approaches established that sequential deep learning could match or exceed rule-based
tools on their native benchmark distributions.
However, they shared important weaknesses: models trained on slices centred on specific
API calls do not generalise to CWEs where the dangerous operation is not API-identifiable,
and the slice extraction step requires language-specific tooling, preventing direct
application to new languages.

\section{Graph-Based Vulnerability Detection}

A conceptually richer representation of code structure is the program graph: abstract syntax
trees (ASTs), control-flow graphs (CFGs), and data-flow graphs (DFGs) capture relationships
between program elements that are invisible in token sequences.

Devign~\citep{zhou2019devign} applied a gated graph neural network to a combined graph
representation --- a joint encoding of the AST, CFG, call graph, and data-dependence graph ---
and achieved state-of-the-art results on a large C/C++ dataset.
Reveal~\citep{chakraborty2022reveal} similarly used a graph neural network over a
code property graph, and showed that graph representations substantially outperform
sequence representations on functions where the vulnerability depends on inter-statement
relationships such as use-after-free and null pointer dereference.

The practical limitation of graph-based approaches is the complexity of their preprocessing
pipelines: extracting program graphs requires a full compiler or parser for each supported
language, and the graph construction is language-specific.
Building a graph-based system that handles both C and Python without separate per-language
toolchains is substantially harder than adapting a text-based model.

\section{Pre-Trained Code Models}

The availability of large corpora of open-source code has made it possible to pre-train
transformer encoders on code at scale, producing representations that capture both syntactic
and functional regularities across millions of functions.

CodeBERT~\citep{feng2020codebert} was the first large pre-trained model for code,
training a RoBERTa~\citep{liu2019roberta} encoder jointly on natural language and code
using masked language modelling and replaced token detection objectives.
CodeBERT's pre-training on a six-language corpus (Python, Java, JavaScript, PHP, Ruby, Go)
gave it representations that transfer across these languages, and fine-tuning on
vulnerability detection tasks has produced competitive results.

GraphCodeBERT~\citep{guo2021graphcodebert} extended CodeBERT by incorporating
data-flow information into the pre-training objective: in addition to masked language
modelling, the model learns to predict which tokens depend on which other tokens through
an auxiliary data-flow alignment task.
This gives GraphCodeBERT partial access to program structure without requiring full graph
extraction, and it has shown improvements over CodeBERT on code understanding tasks including
vulnerability detection.

UniXcoder~\citep{guo2022unixcoder} further extends this line of work with cross-modal
pre-training objectives that jointly train on code tokens, natural language docstrings,
and linearised abstract syntax trees.
This multi-modal training gives UniXcoder richer representations of code structure than
models trained on code text alone, and its multi-language pre-training corpus makes it
a natural choice for cross-language transfer experiments.
UniXcoder is selected as the base model for this work on this basis.

LineVul~\citep{fu2022linevul} fine-tuned CodeBERT specifically for vulnerability detection
and proposed line-level localisation in addition to function-level classification, showing
that the model's attention mechanism can be used to identify which lines within a vulnerable
function are most likely to contain the flaw.

\section{Vulnerability Datasets}

The availability of labelled vulnerability data is a fundamental bottleneck for ML-based
detection.
BigVul~\citep{fan2020bigvul} is a large-scale dataset of C/C++ functions extracted from
open-source projects, labelled by matching commit messages against CVE records.
It contains approximately 188{,}000 functions, of which around 10{,}000 are labelled vulnerable.
The dataset's principal weakness is label noise: commit message matching is imprecise,
and many functions labelled as vulnerable contain multiple unrelated changes from which the
actual vulnerability is difficult to isolate.

DiverseVul~\citep{ni2023diversevul} addresses this by extracting function-level samples
from commits that reference CVE identifiers in a structured way, producing approximately
330{,}000 labelled samples across a broader range of CWE types than BigVul.
DiverseVul includes structured CWE annotations for a substantial fraction of its entries,
making it possible to evaluate per-class recall rather than only overall accuracy.
It is used in this work as the source of real-world evaluation and augmentation data.

A notable gap in both datasets is Python coverage.
BigVul and DiverseVul are almost exclusively C/C++; large-scale function-level Python
vulnerability datasets comparable in size and annotation quality do not currently exist.
This gap motivates the synthetic data approach used in this work.

\section{Synthetic Vulnerability Data}

The use of synthetic data to supplement or replace real-world labelled data has precedent
in several domains.
In the security context, synthetic vulnerability generation has been used to augment
small real-world datasets and to produce balanced class distributions.
Hand-crafted synthetic examples have the advantage of being clean (no label noise),
structurally precise (the dangerous operation is unambiguous), and extensible to any
vulnerability class the designer can characterise.

The principal risk of synthetic training data is distribution shift: a model trained on
idealised, compact, purpose-built examples may fail to recognise the same vulnerability
when it appears in the longer, noisier, more contextually complex functions found in
real codebases.
This synthetic-to-real generalisation gap is a central focus of this dissertation
and is measured quantitatively in Experiment~3.

\section{Cross-Language Transfer}

Multi-language vulnerability detection has received limited attention in the literature.
Most published systems are evaluated on a single language, and cross-language performance
is rarely reported.
An exception is work using code-to-code translation as an intermediate step: functions
in one language are translated to another (using a neural machine translation model),
and the translated function is evaluated by a monolingual classifier.
This approach is impractical at scale and does not generalise to CWEs that are expressed
through language-specific APIs.

Pre-trained multi-language models offer a more natural path to cross-language transfer:
by sharing encoder parameters across languages, the model learns representations that
should generalise across language boundaries for conceptually shared constructs.
However, the extent to which this transfer holds at the CWE level --- and which CWE classes
benefit most --- has not been systematically investigated.
The cross-language ablation study in Chapter~\ref{ch:experiments} provides the first
detailed analysis of this question across a controlled set of shared vulnerability classes.

\section{Calibration of Neural Classifiers}

A well-calibrated classifier is one whose output probabilities correspond to empirical
frequencies: a prediction of 0.8 should be correct approximately 80\% of the time.
\citet{guo2017calibration} showed that modern neural networks are systematically
overconfident --- their confidence scores are higher than their empirical accuracy ---
and proposed temperature scaling as a simple post-hoc calibration method.
Temperature scaling divides the pre-sigmoid logits by a scalar $T > 1$ to flatten the
output distribution, with $T$ optimised to minimise the Expected Calibration Error (ECE)
on a held-out calibration set.

Label smoothing~\citep{muller2019label} is a training-time alternative that directly
regularises overconfidence by replacing hard targets with soft targets:
$\tilde{y} = y(1 - \varepsilon) + \varepsilon/K$, where $K$ is the number of classes
and $\varepsilon$ is the smoothing parameter.
For binary classification this reduces to $\tilde{y} = y(1-\varepsilon) + \varepsilon \cdot 0.5$.
\citet{muller2019label} showed that label smoothing improves calibration across a wide
range of tasks and can substitute for post-hoc temperature scaling.

Both techniques are applied in this work: label smoothing ($\varepsilon = 0.1$) is used
during training, and temperature scaling is applied post-hoc.
The near-unity temperature found by calibration ($T = 0.981$) confirms that label
smoothing has already produced well-calibrated outputs before post-hoc adjustment.

\section{Summary}

The literature establishes that pre-trained code models are effective for within-language
vulnerability detection at function level, that graph-based representations improve
performance on inter-statement vulnerability classes, and that the availability of
labelled real-world data is a fundamental bottleneck.
The contributions of this dissertation --- a multi-language synthetic dataset, a
cross-language ablation study at CWE level, a semantic understanding test suite, and
targeted training interventions --- address gaps that the existing literature has
not filled.
